%!TEX TS-program = xelatex
%!TEX encoding = UTF-8 Unicode
% Deedy CV/Resume Template
% Original author: Debarghya Das (http://debarghyadas.com)
% Repository: https://github.com/deedy/Deedy-Resume
% License: Apache 2.0
%
% IMPORTANT: Compile with XeLaTeX (requires fontspec + Lato/Raleway fonts)
%

\documentclass[]{deedy-resume-openfont}
\usepackage{fancyhdr}

\pagestyle{fancy}
\fancyhf{}

\begin{document}

\lastupdated

%%%%%%%%%%%%%%%%%%%%%%%%%%%%%%%%%%%%%%
%
%     TITLE NAME
%
%%%%%%%%%%%%%%%%%%%%%%%%%%%%%%%%%%%%%%
\namesection{Hyungtae}{Lim}{
  \urlstyle{same}\href{https://limhyungtae.github.io/hyungtae-lim/}{limhyungtae.github.io} |
  \href{mailto:shapelim@mit.edu}{shapelim@mit.edu} |
  \href{https://scholar.google.com/citations?user=S1A3nbIAAAAJ}{Google Scholar}
}

%%%%%%%%%%%%%%%%%%%%%%%%%%%%%%%%%%%%%%
%
%     COLUMN ONE (LEFT)
%
%%%%%%%%%%%%%%%%%%%%%%%%%%%%%%%%%%%%%%
\begin{minipage}[t]{0.33\textwidth}

%%%%%%%%%%%%%%%%%%%%%%%%%%%%%%%%%%%%%%
%     EDUCATION
%%%%%%%%%%%%%%%%%%%%%%%%%%%%%%%%%%%%%%
\section{Education}

\subsection{KAIST}
\descript{Ph.D. in EE \& Robotics}
\location{Mar. 2020 -- Feb. 2023 | Daejeon, Korea}
Advisor: Prof. Hyun Myung \\
Dissertation: Robust LiDAR SLAM \\
for Autonomous Vehicles
\sectionsep

\subsection{KAIST}
\descript{M.S. in EE \& Robotics}
\location{Mar. 2018 -- Feb. 2020 | Daejeon, Korea}
\sectionsep

\subsection{KAIST}
\descript{B.S. in Mechanical Engineering}
\location{Mar. 2013 -- Aug. 2018 | Daejeon, Korea}
Dean's List (Fall 2015)
\sectionsep

%%%%%%%%%%%%%%%%%%%%%%%%%%%%%%%%%%%%%%
%     LINKS
%%%%%%%%%%%%%%%%%%%%%%%%%%%%%%%%%%%%%%
\section{Links}
Website:// \href{https://limhyungtae.github.io/hyungtae-lim/}{\bf limhyungtae.github.io} \\
GitHub:// \href{https://github.com/limhyungtae}{\bf limhyungtae} \\
LinkedIn:// \href{https://www.linkedin.com/in/hyungtae-lim}{\bf hyungtae-lim} \\
Scholar:// \href{https://scholar.google.com/citations?user=S1A3nbIAAAAJ}{\bf Hyungtae Lim}
\sectionsep

%%%%%%%%%%%%%%%%%%%%%%%%%%%%%%%%%%%%%%
%     SKILLS
%%%%%%%%%%%%%%%%%%%%%%%%%%%%%%%%%%%%%%
\section{Skills}
\subsection{Programming}
\location{Proficient:}
C++ \textbullet{} Python \textbullet{} ROS/ROS2 \\
\location{Experienced:}
PyTorch \textbullet{} CUDA \textbullet{} CMake \\
OpenCV \textbullet{} \LaTeX
\sectionsep

\subsection{Research}
SLAM \textbullet{} 3D LiDAR \\
Point Cloud Registration \\
Ground Segmentation \\
Sensor Fusion \textbullet{} Deep Learning \\
Multi-Robot Mapping
\sectionsep

\subsection{Languages}
Korean (Native) \textbullet{} English (Fluent)
\sectionsep

\end{minipage}
\hfill
%%%%%%%%%%%%%%%%%%%%%%%%%%%%%%%%%%%%%%
%
%     COLUMN TWO (RIGHT)
%
%%%%%%%%%%%%%%%%%%%%%%%%%%%%%%%%%%%%%%
\begin{minipage}[t]{0.66\textwidth}

%%%%%%%%%%%%%%%%%%%%%%%%%%%%%%%%%%%%%%
%     EXPERIENCE
%%%%%%%%%%%%%%%%%%%%%%%%%%%%%%%%%%%%%%
\section{Experience}

\runsubsection{SPARK Lab, MIT}
\descript{| Postdoctoral Associate}
\location{Apr. 2024 -- Present | Cambridge, MA}
\vspace{\topsep}
\begin{tightemize}
\item Advised by Prof. Luca Carlone (LIDS, MIT).
\item Research on multi-robot/-session autonomous mapping in urban environments.
\item Guide graduate and undergraduate students in research papers and theses.
\end{tightemize}
\sectionsep

\runsubsection{Urban Robotics Lab, KAIST}
\descript{| Postdoctoral Fellow}
\location{Mar. 2023 -- Mar. 2024 | Daejeon, Korea}
\begin{tightemize}
\item Research on LiDAR and vision-based localization and mapping.
\item Developed deep learning models for robust robot perception.
\end{tightemize}
\sectionsep

\runsubsection{StachnissLab, Univ. Bonn}
\descript{| Visiting Scholar}
\location{Nov. 2022 -- Feb. 2023 | Bonn, Germany}
\begin{tightemize}
\item Collaborated with Prof. Cyrill Stachniss on static map building.
\item Paper accepted at RSS 2023: \textit{ERASOR2}.
\end{tightemize}
\sectionsep

\runsubsection{NAVER LABS}
\descript{| Research Intern}
\location{Apr. 2022 -- Sep. 2022 | Gyeonggi-Do, Korea}
\begin{tightemize}
\item Worked on robust global 3D point cloud registration.
\item Paper accepted at IEEE ICRA 2023: \textit{ORORA}.
\end{tightemize}
\sectionsep

%%%%%%%%%%%%%%%%%%%%%%%%%%%%%%%%%%%%%%
%     SELECTED PUBLICATIONS
%%%%%%%%%%%%%%%%%%%%%%%%%%%%%%%%%%%%%%
\section{Selected Publications}
\begin{tightemize}
\item \textbf{H. Lim} et al., ``KISS-Matcher: Fast and Robust Point Cloud Registration Revisited,'' \textit{IEEE ICRA}, 2025.
\item D. Maggio, \textbf{H. Lim}, and L. Carlone, ``VGGT-SLAM: Dense RGB SLAM Optimized on the SL(4) Manifold,'' \textit{NeurIPS}, 2025.
\item M. Seo, \textbf{H. Lim} et al., ``BUFFER-X: Towards Zero-Shot Point Cloud Registration,'' \textit{IEEE/CVF ICCV}, 2025.
\item \textbf{H. Lim} et al., ``ERASOR2: Instance-Aware Robust 3D Mapping of the Static World,'' \textit{RSS}, 2023.
\item \textbf{H. Lim} et al., ``Quatro++: Robust Global Registration Exploiting Ground Segmentation,'' \textit{IJRR}, 2023.
\item \textbf{H. Lim} et al., ``Patchwork: Concentric Zone-Based Ground Segmentation,'' \textit{IEEE RA-L w/ IROS}, 2021.
\item \textbf{H. Lim} et al., ``ERASOR: Egocentric Ratio-Based Dynamic Object Removal,'' \textit{IEEE RA-L w/ ICRA}, 2021.
\end{tightemize}
\sectionsep

%%%%%%%%%%%%%%%%%%%%%%%%%%%%%%%%%%%%%%
%     AWARDS
%%%%%%%%%%%%%%%%%%%%%%%%%%%%%%%%%%%%%%
\section{Awards}
\begin{tabular}{rll}
2025 & Outstanding & IEEE ICRA'25 Reviewer Award (5/7,641) \\
2024 & Pioneer     & RSS Pioneers 2024 (30/202 selected) \\
2024 & 1\textsuperscript{st} Prize & HILTI SLAM Challenge, IEEE ICRA'24 \\
2023 & Best Paper  & IEEE RA-L (5 of 1,100 papers) \\
2023 & 1\textsuperscript{st} Prize & HILTI SLAM Challenge, IEEE ICRA'23 \\
2023 & Innovation  & CES'23 Innovation Award \\
2022 & 2\textsuperscript{nd} Cash  & HILTI SLAM Challenge, IEEE ICRA'22 \\
2020 & Best Paper  & Student Best Paper Award, ICCAS \\
\end{tabular}
\sectionsep

\end{minipage}
\end{document}
